% ==================================
% 1. MACROS SEMÁNTICAS HISTÓRICAS
% ==================================
% Griegos
\providecommand{\Episteme}{\textit{episteme}}
\providecommand{\Doxa}{\textit{doxa}}
\providecommand{\Logos}{\textit{logos}}
\providecommand{\Aisthesis}{\textit{aisthesis}}
\providecommand{\Nous}{\textit{nous}}
\providecommand{\Ousia}{\textit{ousia}}
\providecommand{\Psuche}{\textit{psuché}}
\providecommand{\Epagoge}{\textit{epagoge}}

% Medievales
\providecommand{\SiFallor}{\textit{Si fallor, sum}}
\providecommand{\Suppositio}{\textit{suppositio}}
\providecommand{\NotitiaIntuitiva}{\textit{notitia intuitiva}}
\providecommand{\NotitiaAbstractiva}{\textit{notitia abstractiva}}

% Modernidad
\providecommand{\Cogito}{\textit{cogito}}
\providecommand{\ResCogitans}{\textit{res cogitans}}
\providecommand{\ResExtensa}{\textit{res extensa}}
\providecommand{\DudaMetodica}{duda metódica}
\providecommand{\GenioMaligno}{genio maligno}
\providecommand{\IdeasInnatas}{ideas innatas}
\providecommand{\ClaridadDistincion}{claridad y distinción}
\providecommand{\Fundacionalismo}{fundacionalismo}

% ==================================
% 2. MACROS SEMÁNTICAS TEMÁTICAS
% ==================================
% Ontología Social y Lenguaje (Searle / Austin)
\providecommand{\HechoBruto}{\textit{hecho bruto}}
\providecommand{\HechosBrutos}{\textit{hechos brutos}}
\providecommand{\HechoInstitucional}{\textit{hecho institucional}}
\providecommand{\HechosInstitucionales}{\textit{hechos institucionales}}
\providecommand{\FuncionEstatus}{\textit{función de estatus}}
\providecommand{\ReglaRegulativa}{regla regulativa}
\providecommand{\ReglaConstitutiva}{regla constitutiva}
\providecommand{\ActoDeclarativo}{acto de habla declarativo}
\providecommand{\Trasfondo}{\textit{Trasfondo}}
\providecommand{\EfectoTinkerbell}{Efecto Tinkerbell}

% Poder, Cultura y Materialismo Histórico (Marx / Bourdieu / Foucault)
\providecommand{\BaseMaterial}{base material}
\providecommand{\Superestructura}{superestructura}
\providecommand{\Proletarizacion}{proletarización cognitiva}
\providecommand{\CajaNegra}{\textbf{caja negra}}
\providecommand{\SavoirFaire}{\textit{savoir-faire}}
\providecommand{\SavoirVivre}{\textit{savoir-vivre}}
\providecommand{\Habitus}{\textit{habitus}}
\providecommand{\RegimenVerdad}{régimen de verdad}
\providecommand{\ViolenciaSimbolica}{violencia simbólica}
\providecommand{\FetichismoMercancia}{fetichismo de la mercancía}

% Antropología Filosófica (Cassirer / Blumenberg / Nietzsche)
\providecommand{\AnimalSimbolico}{\textit{animal simbólico}}
\providecommand{\AbsolutismoRealidad}{absolutismo de la realidad}
\providecommand{\FormaSimbolica}{\textit{forma simbólica}}
\providecommand{\FormasSimbolicas}{\textit{formas simbólicas}}
\providecommand{\ParsProToto}{\textit{pars pro toto}}
\providecommand{\EficaciaSimbolica}{eficacia simbólica}

% ==================================
% 3. MARCOS METODOLÓGICOS Y ANALÍTICOS
% ==================================
\providecommand{\MetodoACP}{\textbf{Método ACP}}
\providecommand{\MetodoMCZ}{\textbf{MCZ}} % Método Cornell-Zettelkasten
\providecommand{\ModeloTresC}{\textbf{M-3C}} % Modelo Tripartito
\providecommand{\MTAM}{\textbf{M-TAM}} % Modelo Tetrádico para el Análisis Metafísico

% ==================================
% 4. MACROS LÓGICAS Y MATEMÁTICAS
% ==================================
% Lógica Proposicional y Conjuntos
\providecommand{\K}{\mathsf{K}}
\providecommand{\V}{\mathsf{V}}
\providecommand{\Implica}{\Rightarrow}
\providecommand{\Sii}{\Leftrightarrow}

% Laboratorio 1: Antropogénesis y Epifilogénesis
\providecommand{\Org}[1]{O_{#1}}        % Organismo en t
\providecommand{\Tec}[1]{T_{#1}}        % Técnica en t
\providecommand{\Med}{M}                % Medio ambiente
\providecommand{\Hum}{H}                % Humano como suma
\providecommand{\mGen}{m_{\text{gen}}}  % Memoria Genética
\providecommand{\mEpi}{m_{\text{epi}}}  % Memoria Epigenética
\providecommand{\mTec}{m_{\text{tec}}}  % Memoria Tecnológica

% Laboratorio 2: Termodinámica del Poder Institucional
\providecommand{\VarX}{X}               % Soporte físico (Hecho bruto)
\providecommand{\VarY}{Y}               % Función de estatus
\providecommand{\VarC}{C}               % Contexto normativo
\providecommand{\Legit}{L(t)}           % Legitimidad 
\providecommand{\Viol}{V(t)}            % Violencia 
\providecommand{\ConstK}{K}             % Constante de estabilidad
\providecommand{\Edisp}{E_{\text{disponible}}} % Energía base

% Laboratorio 3: Termodinámica Psíquica del Símbolo
\providecommand{\VarR}{R}               % Lo Real / Caos sensorial
\providecommand{\VarS}{S}               % Sistema Simbólico
\providecommand{\VarP}{P}               % Percepción Consciente
\providecommand{\VarA}{A(t)}            % Angustia / Entropía psíquica
\providecommand{\Varx}{x}               % Elemento singular (la parte)
\providecommand{\VarXTodo}{X}           % El Todo / Estructura general

% M-TAM Formalización
\providecommand{\DominioW}{W}           % Mundo
\providecommand{\DominioC}{C}           % Conciencia
\providecommand{\DominioP}{P}           % Fenómenos
\providecommand{\DominioS}{S}           % Símbolos

% ==================================
% 5. MACROS SEMÁNTICAS (Sesión 4: El Gran Cisma)
% ==================================
% Ferdinand Tönnies (Comunidad vs. Sociedad)
\providecommand{\Gemeinschaft}{\textit{Gemeinschaft}}
\providecommand{\Gesellschaft}{\textit{Gesellschaft}}
\providecommand{\Wesenwille}{\textit{Wesenwille}}
\providecommand{\Kurwille}{\textit{Kürwille}}
\providecommand{\AcumulacionOriginaria}{acumulación originaria}
\providecommand{\Cercamientos}{\textit{Enclosure Acts}}

% Émile Durkheim & Georg Simmel (Patología Urbana)
\providecommand{\Anomia}{\textit{anomia}}
\providecommand{\ActitudBlase}{\textit{actitud blasé}}
\providecommand{\Blasiertheit}{\textit{Blasiertheit}}

% Max Weber (Racionalización)
\providecommand{\Desencantamiento}{\textit{Entzauberung der Welt}}
\providecommand{\RacionalidadInstrumental}{\textit{Zweckrationalität}}
\providecommand{\JaulaDeHierro}{\textit{stahlhartes Gehäuse}}

% ==================================
% 6. MACROS FORMALES (Sesión 4: Termodinámica del Vínculo Social)
% ==================================
\providecommand{\VarLib}{L(t)}          % Libertad Individual (Movilidad, anonimato)
\providecommand{\VarSeg}{S(t)}          % Seguridad Ontológica (Pertenencia)
\providecommand{\ConstCoh}{k}           % Constante de cohesión social
\providecommand{\VarNorm}{N(t)}         % Estructura normativa / Nivel de contención
\providecommand{\VarDes}{D(t)}          % Nivel de Deseo infinito
\providecommand{\VarPos}{P(t)}          % Posibilidad material real
\providecommand{\VarAno}{A(t)}          % Anomia estructural
\providecommand{\VarRac}{R(t)}          % Nivel de Racionalización/Burocracia
\providecommand{\VarEfi}{E(t)}          % Eficiencia técnica
\providecommand{\VarMen}{M(t)}          % Sentido / Espíritu (Meaning)
\providecommand{\VarPsi}{\Psi(t)}       % Costo Psíquico

% ==================================
% ENTORNOS SEMÁNTICOS
% ==================================
\newtcolorbox{notaautor}{
  colback=white,
  colframe=black,
  title=Nota de investigación,
  fonttitle=\bfseries,
  sharp corners,
  breakable
}

\newtcolorbox{objetivo}{
  colback=white,
  colframe=black,
  title=Objetivos,
  fonttitle=\bfseries,
  sharp corners,
  breakable
}

\newtcolorbox{actividad}{
  colback=white,
  colframe=black,
  title=Actividad,
  fonttitle=\bfseries,
  sharp corners,
  breakable
}

\newtcolorbox{argumento}{
  colback=white,
  colframe=black,
  title=Estructura argumental,
  fonttitle=\bfseries,
  sharp corners,
  breakable
}

\newtcolorbox{cita}{
  colback=white,
  colframe=black,
  sharp corners,
  breakable
}

\newtcolorbox{problema}{
  colback=white,
  colframe=black,
  title=Problema,
  fonttitle=\bfseries,
  sharp corners,
  breakable
}

\newtcolorbox{definicion}{
  colback=white,
  colframe=black,
  title=Definición,
  fonttitle=\bfseries,
  sharp corners,
  breakable
}

\newtcolorbox{objecion}{
  colback=white,
  colframe=black,
  title=Objeción,
  fonttitle=\bfseries,
  sharp corners,
  breakable
}

\newtcolorbox{ejemplo}{
  colback=white,
  colframe=black,
  title=Ejemplo,
  fonttitle=\bfseries,
  sharp corners,
  breakable
}

% ==================================
% NOTAS Y COMANDOS DE TRABAJO
% ==================================
\newcommand{\margennota}[1]{\marginnote{\footnotesize #1}}
\newcommand{\TODOCITA}[1]{\todo[inline]{Insertar cita: #1}}
\newcommand{\TODOARG}[1]{\todo[inline]{Revisar argumento: #1}}
\newcommand{\TODOTEX}[1]{\todo[inline]{Revisar redacción: #1}}

\ifmanuscrito
  \newcommand{\notaviva}[1]{\begin{notaautor}#1\end{notaautor}}
\else
  \newcommand{\notaviva}[1]{}
\fi
